\section{Introduction}
The mathematical problem solving capabilities of Large Language Models (LLMs) have greatly improved in recent years. Much of the credit for LLMs successes in this field involve the use of Chain-of-Thought, prompting the model to output a step-by-step solution to the problem, which often elicits a virtuous trajectory where explicit solution steps might reinforce the direction of the correct answer \cite{chain-of-thought}. In the same way, however, incorrect solution steps can lead the model astray and much effort has been spent in verifying individual reasoning steps, starting with big annotated datasets used to fine-tune external judging LLMs \cite{lightman2023lets}. More recent literature, however, has attempted to use more transparent and deterministic verifiers by using LLMs to translate the natural language solution in a formal language such as Lean \cite{safe2025, stepwise2025}. Even though such languages are usually considered a natural choice to deterministically verify mathematical statements, LLMs had considerably less exposure to such languages during training, as opposed to popular general programming languages like Python. Even models specifically fine-tuned for autoformalisation are far from perfect, with out-of-distribution performance deterioration \cite{putnambench2024} and sycophancy \cite{cornish2026faithformbenchbenchmarkingfaithfulnessmathematical} among the problems in this context, which are particularly relevant when verifying simple algebraic reasoning steps which are typically far from how these languages are used and, therefore, relatively under-represented in the training data. To tackle this problem, we propose a hybrid verification approach where a router is tasked with routing each individual steps to one of two verification pipelines: for simple computation-heavy steps, Python is a natural choice as a very popular coding language for which LLMs have shown remarkable capabilities, while for steps requiring more complicated computation and higher-order logic we have chosen Lean as a formal language which has recently gained traction in the automatic theorem proving community. We show how a simple, small, open-source LLM can suffice to act as such router, being competitive with proprietary, more capable choices, while fine-tuning a BERT model \cite{devlin-etal-2019-bert} on a training corpus for the task also achieve competitive results at an even lower price. Finally, in each case we show how combining Python and Lean does lead to a better verification performance overall and we show what linguistic expression is specifically associated with Python or Lean by analysing the predictions of the fine-tuned BERT router.
